% Generated by `taskgen course examples` from number_theory.fermat_test@carmichael (seed 16).
% Do not edit: change examples.yaml and regenerate.
\begin{example}[тест Ферма для числа Кармайкла]\label{ex:fermat-test-carmichael}
Выполнить тест Ферма для числа $N = 1729$ с основаниями $a = 2, 3, 5$ (по порядку; остановиться на первом свидетеле составности). Каков вердикт теста? Затем проверить его, разложив $N$ на множители, и, если $N$ составное, проверить для $N$ критерий Корсельта: является ли $N$ числом Кармайкла?

\par\noindent\textit{Решение.} Для простого $N$ и $\gcd(a, N) = 1$ по малой теореме Ферма $a^{N-1} \equiv 1 \pmod N$; если $\gcd(a, N) > 1$ или $a^{N-1} \not\equiv 1$, то $a$ --- свидетель составности. $N - 1 = 1728 = (11011000000)_2$. $\gcd(2, 1729) = 1$. $\gcd(3, 1729) = 1$. $\gcd(5, 1729) = 1$. Возведение в степень $N - 1$ слева направо ($r \gets r^2$, при бите $1$ ещё $r \gets r a$; $e$ --- уже обработанная часть показателя):\[\begin{array}{c|r|r|r|r}\text{бит} & e & 2^{e} \bmod N & 3^{e} \bmod N & 5^{e} \bmod N \\ \hline 1 & 1 & 2 & 3 & 5 \\ 1 & 3 & 8 & 27 & 125 \\ 0 & 6 & 64 & 729 & 64 \\ 1 & 13 & 1276 & 185 & 1461 \\ 1 & 27 & 645 & 664 & 1217 \\ 0 & 54 & 1065 & 1 & 1065 \\ 0 & 108 & 1 & 1 & 1 \\ 0 & 216 & 1 & 1 & 1 \\ 0 & 432 & 1 & 1 & 1 \\ 0 & 864 & 1 & 1 & 1 \\ 0 & 1728 & 1 & 1 & 1 \\ \end{array}\]$2^{1728} \equiv 1$ --- раунд пройден; $3^{1728} \equiv 1$ --- раунд пройден; $5^{1728} \equiv 1$ --- раунд пройден.\par\noindent \textbf{Вердикт:} $N$ --- вероятно простое по основаниям $2, 3, 5$.\par\noindent \textbf{Проверка.} $N = 7 \cdot 13 \cdot 19$ --- составное: тест Ферма \emph{ошибся}. Критерий Корсельта: $N$ --- число Кармайкла $\iff$ $N$ свободно от квадратов и $p - 1 \mid N - 1$ для всех $p \mid N$. $1728 \bmod 6 = 0$, $1728 \bmod 12 = 0$, $1728 \bmod 18 = 0$ --- все делимости выполнены: $N$ --- \textbf{число Кармайкла}, $a^{N-1} \equiv 1$ для \emph{всех} $a$, взаимно простых с $N$; выявить составность тестом Ферма может только основание, имеющее общий делитель с $N$.
\end{example}
